\section{Curation Framework}
\label{sec:framework}

Investigating the boundaries of digital library curation requires focusing on what curators can reliably assert when open \llm artifact metadata is scattered across model hubs and scholarly literature. Rather than attempting to reconstruct the exhaustive training history of an artifact or executing the model to verify performance, our evaluation measures public-record sufficiency: whether visible traces support an auditable curation record. Public traces include model repository pages, model cards, links or textual mentions within papers, upstream model fields, family labels, code repository links, and formal release statements. Retaining the explicit source, specific role, and evidentiary strength of these traces allows digital libraries to convert surface-level visibility into verifiable curation evidence, following broader work on FAIR artifacts, software citation, research-object packaging, and context-sensitive scholarly links~\cite{wilkinson2016fair,smith2016softwarecitation,soilandreyes2022rocrate,salsabil2025context}.

Synthesizing a single paper record with its corresponding model-hub repositories allows curators to construct an evidence-backed curation record. Such a record addresses four core dimensions: what the artifact is, how it connects to specific scholarship, what upstream dependencies or derivation relations are publicly supported, and which release assets are visible for preservation triage. These dimensions draw from model documentation, provenance, and reproducibility research, where artifact identity, context, derivation, and inspectable release materials are treated as separate metadata requirements rather than a single availability flag~\cite{mitchell2019modelcards,gebru2021datasheets,moreau2013prov,pineau2021reproducibility}. Incorporating an integrated curatability status explicitly marks whether these distinct evidentiary fields co-occur within a single record. Shifting the assessment framework from binary accessibility metrics, such as link presence or download capability, allows digital libraries to evaluate record sufficiency rather than mere public visibility.

We therefore evaluate five curation functions. Artifact identity asks whether public records distinguish base checkpoints, fine-tunes, adapters, merges, quantized packages, distilled students, and deployment artifacts. Scholarly linkage asks whether the artifact can be connected to papers, repositories, owners, dates, licenses, and task contexts. Provenance sufficiency asks whether upstream evidence is explicit or canonical, inventory-backed, contextual, or missing, following provenance research that separates traceable derivation relations from weaker contextual evidence~\cite{moreau2013prov,davidson2008provenance,herschel2017survey}. Release-package readiness asks whether visible records expose code, scripts, recipes, checkpoints, adapters, data, or evaluation assets rather than only release intentions, reflecting reproducibility work showing that access statements alone do not establish practical reuse or rerun readiness~\cite{dodge2019show,pineau2021reproducibility,akella2023effort}. Integrated curatability asks whether these fields co-occur in one auditable record.

Two operational boundaries define the scope of this framework. First, the study measures curation claims supported strictly by public records rather than trying to map the absolute, exhaustive genealogy of an artifact. True model lineage frequently depends on private training logs, local commit histories, unreleased intermediate checkpoints, or implicit developer knowledge, forcing digital libraries to rely on visible records to establish auditable metadata. This boundary is especially important for open-model reuse, where public records may expose family names, model-card fields, or downstream adaptation signals without fully specifying the underlying dependency relation~\cite{jiang2023ptmreuse,taraghi2024modelreuse,white2024mof}. Second, evidence must be tiered rather than treating all public signals as equivalent. Model-card fields, paper-side links, family mentions, and formal release statements serve fundamentally different curation purposes and cannot be used interchangeably~\cite{mitchell2019modelcards,pepe2024hfdocs,kadasi2025modelhubs,salsabil2025context}. Preserving the alignment among evidence sources, audit checks, and target uses guides the empirical measurement design across five core parameters: identity, linkage, provenance, release packages, and integrated curatability.
